\documentclass[a4paper, 12pt]{article}

\usepackage{utils}

\renewcommand*{\today}{25 February 2026}

\begin{document}

\hotbox{Analyse 4}{CM 6}{\today}

\section{Approximation de Weierstrass}

\subsection{Polynômes de Bernstein}

\begin{definition}
    Soit $n \in \N$ et $f \in \mathcal{C}^0([0, 1], \R)$. On note $B_n(\cdot, f)$ le
    n-ième polynôme de Bernstein de $f$ défini par
    $$
    B_n(X, f) = \sum_{k=0}^n f\left(\frac{k}{n}\right) \binom{n}{k} X^k (1-X)^{n-k}
    $$ 
\end{definition}

\begin{exemple}
    On prend $B_n(X, x \mapsto 1) = \sum\limits_{k=0}^n 1 \times  \binom{n}{k} X^k (1-X)^{n-k} = (X + (1-X))^n = 1^n = 1$.
\end{exemple}

\begin{exemple}
    On prend 
    \begin{align*}  
        B_n(X, x \mapsto x) &= \sum\limits_{k=0}^n \frac{k}{n} \times  \binom{n}{k} X^k (1-X)^{n-k} \\
        &= \sum\limits_{k=0}^n \frac{k}{n} \dfrac{n!}{k!(n-k)!} X^k (1-X)^{n-k} \\
        &= \sum\limits_{k=0}^n \frac{(n-1)!}{(k-1)!((n-1)-(k-1))!} X^k (1-X)^{(n-1)-(k-1)} \\
        &= x \sum\limits_{k'=0}^{n-1} \binom{n-1}{k'} X^{k'} (1-X)^{(n-1)-(k')} \quad car k < 0 \implies \binom{n}{k} = 0 \\
        &= x \sum\limits_{k=0}^{n-1} \binom{n-1}{k} X^k (1-X)^{(n-1)-k} \\
        &= x (x + (1-x))^{n-1} \\
        &= x
    \end{align*}
\end{exemple}

\begin{exemple}
    On utilise
    \begin{align*}
        B_n(X, x \mapsto x^2) &= \sum\limits_{k=0}^n \left( \frac{k}{n} \right)^2 \binom{n}{k} X^k (1 - X)^{n-k} \\
    \end{align*}
    \begin{hotwarn}
        À faire
    \end{hotwarn}
\end{exemple}

\begin{lemme}{}{}
    $$
    \sum_{k=0}^n (k - nX)^2 \binom{n}{k} X^k (1-X)^{n-k} = nX(1-X)
    $$
\end{lemme}

\subsection{Théorème de Weierstrass}

\begin{definition}
    Soient $(X, d_X)$ et $(Y, d_Y)$ deux espaces métriques. Soit $f: X \to Y$ et $\delta > 0$.

    On appelle module de continuité de $f$, à l'échelle $\delta$, noté $\omega(f, \delta)$ la quantité
    $$
    \omega(f, \delta) = \sup_{d_X(x, y) \leq \delta} d_Y(f(x), f(y))
    $$
\end{definition}

\begin{remarque}
    \begin{align*}
        \lim_{\delta \to 0} \omega(f, \delta) = 0 &\iff \forall \varepsilon > 0, \exists \delta > 0, |\omega(f, \delta) - 0| < \varepsilon \\
        &\iff \forall \varepsilon > 0, \exists \delta > 0, d_X(x, y) \leq \delta \implies d_Y(f(x), f(y)) < \varepsilon \\
        &\iff \forall \varepsilon > 0, \exists \delta > 0, \forall x, y \in I, |x - y| \leq \delta \implies |f(x) - f(y)| < \varepsilon \\
        &\iff f \text{ est uniformément continue sur I}
    \end{align*}
\end{remarque}

\begin{theoreme}{Weierstrass}{}
    Soit $f: \mathcal{C}^0([0, 1], \R)$. Alors
    $B_n(\cdot, f)$ converge uniformément vers $f$ quand $n \to +\infty$ sur $[0, 1]$.
\end{theoreme}

\begin{demonstration}
    Soit $\varepsilon > 0$, $\delta > 0$ et $x \in [0, 1]$ alors on a
    \begin{align*}
        |B_n(x, f) - f(x)| &= \left| \sum_{k=0}^n f\left( \frac{k}{n} \right) \binom{n}{k} x^k (1-x)^{n-k} - f(x) \sum_{k=0}^n \binom{n}{k} x^k (1-x)^{n-k} \right| \\
        &= \left| \sum_{k=0}^n \binom{n}{k} x^k (1-x)^{n-k} \left( f\left( \frac{k}{n} \right) - f(x) \right) \right| \\
        &\leq \sum_{k=0}^n \binom{n}{k} x^k (1-x)^{n-k} \left| f\left( \frac{k}{n} \right) - f(x) \right| \\
        &\leq \sum_{|\frac{k}{n} - x| < \delta} \binom{n}{k} x^k (1-x)^{n-k} \left| f\left( \frac{k}{n} \right) - f(x) \right|\\
        &\quad + \sum_{|\frac{k}{n} - x| \geq \delta} \binom{n}{k} x^k (1-x)^{n-k} \left| f\left( \frac{k}{n} \right) - f(x) \right| \\
    \end{align*}
    On a $\forall x \in [0, 1], \sup_{|\frac{k}{n} - x| < \delta} \left| f\left( \frac{k}{n} \right) - f(x) \right| = \omega(f, \delta)$ donc le premier terme c'est ok.

    Pour le deuxième terme, on va majorer $|f(\frac{k}{n}) - f(x)|$ par $2 \norm{f}_{\infty, [0, 1]}$ et $f$ est continue sur $[0, 1]$ donc elle est fermée bornée (théorème de Heine) alors
    $f$ est uniformément continue sur $[0, 1]$.

    On a pour la somme
    $$
    \left|\dfrac{k}{n} - x \right| \iff \left( \dfrac{k}{n} - x \right)^2 \geq \delta^2
    $$
    donc $\dfrac{\left(\frac{k}{n} - x\right)^2}{\delta^2} \geq 1$

    Ainsi
    \begin{align*}
        I &= \sum_{\left|\frac{k}{n} - x \right| > \delta^2} \binom{n}{k} x^k (1-x)^{n-k} \left| f\left( \frac{k}{n} \right) - f(x) \right| \\
        &\leq \sum_{\left|\frac{k}{n} - x \right| > \delta^2} \binom{n}{k} x^k (1-x)^{n-k} 2 \norm{f}_{\infty, [0, 1]} \\
    \end{align*}
    \begin{hotwarn}
        À faire
    \end{hotwarn}
\end{demonstration}

\end{document}